\section{What changes, and what remains true}
\label{sec:change}

\subsection{Enlarging a test is not refuting its restricted version}

There is an easy but misleading way to describe the example: homotopy theory
reveals that the set-valued resolution was false. It does not. The
set-valued assertion remains true for every set. What changes is the range
of objects over which sufficiency is quantified.

\begin{proposition}[Persistence of the earlier tests]\label{prop:persistence}
In the setting of \cref{lem:criterion}, let $0\leq k\leq\ell\leq\infty$.
For every $A\in\mathcal V_k$, the resolution comparison is the same when
computed in $\mathcal V_k$ or in $\mathcal V_\ell$. Hence an
$\ell$-resolver is a $k$-resolver, while a $k$-resolver can fail to be an
$\ell$-resolver only on additional test objects.
\end{proposition}

\begin{proof}
The inclusion $\mathcal V_k\hookrightarrow\mathcal V_\ell$ is fully
faithful and preserves limits. Left extension from $D$ is constant in both
settings. Thus the pointwise right Kan extension and its unit agree on
$A\in\mathcal V_k$. An equivalence for all objects of the larger
category is, in particular, one for all objects of the smaller category.
\end{proof}

This proposition locates the result exactly. It does not exhibit a
contradiction between two answers to the same fully specified mathematical
question. It exhibits different answers to related questions whose
coefficient parameter was easily left implicit. The philosophical problem is
to determine when that parameter may legitimately remain fixed and when
the interpretation itself requires it to vary.

In particular, our main family is stronger as a test than the familiar
observation that a homotopy colimit can differ from an ordinary colimit.
That observation is relevant background \cite{Dugger}, and an elementary
example is recorded in the appendix. Here the source construction is
unchanged even on nondiscrete coefficients: $\Lan_D A$ is constant.
The defect lies in its reconstruction through a higher right Kan
extension. We have isolated the test of resolution from a possible change
in the generation of the object being tested.

More generally, let $i:\mathcal V\hookrightarrow\mathcal W$ be
fully faithful and preserve the relevant limits. If the canonical
comparison
\[
i\Lan_D^{\mathcal V} A\longrightarrow\Lan_D^{\mathcal W}(iA)
\]
is an equivalence for every old input $A$, then the generated resolution
units agree on those inputs. For finite $C$, this implies
$\Aup_{\mathcal V}(D)\subseteq\Aup_{\mathcal W}(D)$: a resolver in
$\mathcal W$ passes every old comparison, and hence contains the least
one in $\mathcal V$. The initial-singleton case satisfies this hypothesis,
as does the nine-object example because its generation uses coproducts.
The hypothesis is sufficient, and the next result shows why it cannot
simply be omitted.

\subsection{Changing generation can shrink a least resolver}
\label{sec:reversal}

Preservation of left generation is a substantive hypothesis in the
comparison of coefficients. Without it, the least resolving subset can
become smaller when the coefficient category becomes larger.

\begin{proposition}[Reversal when generation changes]
\label{prop:reversal}
Let $J$ be a nonempty finite poset, and put
\[
C=\{x\}\star J\star\{t\},\qquad D=J,\qquad U=J\star\{t\}.
\]
For finite $k\geq0$ let $T_k=\tau_{\leq k}\B J$, and put
$T_\infty=\B J$. Then
\[
\Aup_k(D)=
\begin{cases}
C,&T_k\simeq *,\\
U,&T_k\not\simeq *.
\end{cases}
\]
In particular, if $J=Q_{d+1}$ is the ordinal sum of $d+1$
two-element antichains, where $d\geq1$, then $|C|=2d+4$ and
\[
\Aup_k(D)=
\begin{cases}
C,&0\leq k<d,\\
C\setminus\{x\},&d\leq k\leq\infty.
\end{cases}
\]
\end{proposition}

\begin{proof}
For every $A:J\to\mathcal V_k$, its left extension $X=\Lan_D A$
has $X(x)=\varnothing$, restricts to $A$ on $J$, and has
$X(t)=\colim_J A$ with the canonical cocone. The terminal position
$t$ is mandatory at every coefficient level: taking the empty input
diagram gives the matching map $\varnothing\to *$ at $t$. The
positions in $D$ are mandatory by definition. It therefore remains
only to decide whether $x$ is mandatory.

Take the constant terminal input, and write
$F=(\Lan_D1)|_U$. Its restriction to $J$ is terminal and
$F(t)=T_k$. The canonical cocone represents the truncation map
$\eta:\B J\to T_k$, or the identity for $k=\infty$. A cone over
$F$ is a cone over its restriction to $J$, a point of $T_k$, and
compatibility with that cocone. Consequently
\[
\limop_U F\simeq
T_k\times_{\Map(\B J,T_k)}\{\eta\},
\]
where $T_k\to\Map(\B J,T_k)$ sends a point to the corresponding
constant map. By the universal property of truncation,
$\Map(\B J,T_k)\simeq\Map(T_k,T_k)$; under this equivalence,
$\eta$ is the identity and constant maps remain constant.
The displayed homotopy pullback is therefore nonempty exactly when
the identity of $T_k$ is homotopic to a constant, equivalently when
$T_k$ is contractible.

If $T_k$ is noncontractible, this limit is empty. For arbitrary
$A$, the unique map $A\to1$ induces
\[
\limop_U(\Lan_D A)|_U\longrightarrow\limop_U F=\varnothing.
\]
An object of spaces or truncated spaces admitting a map to the
empty object is itself empty. Thus the matching map at $x$ is
$\varnothing\to\varnothing$ for every input, and $U$ resolves $D$.
Together with the mandatory positions already identified, this
proves leastness.

If $T_k$ is contractible, $F$ is pointwise terminal on $U$, so its
limit is terminal. The matching map at $x$ for the terminal input
is then $\varnothing\to *$. Thus $x$ is also mandatory, and the
least resolver is $C$.

Finally, $\B Q_{d+1}\simeq S^d$. Its $k$-truncation is
contractible exactly for $0\leq k<d$, and is noncontractible for
finite $k\geq d$ and for unrestricted spaces. This proves the
last formula.
\end{proof}

For $d=1$, the least resolver has six objects for sets and five for
groupoids and all higher coefficients. This reverses the direction of
the main family's change. The reason is explicit: the terminal input
is generated with value $*$ at $t$ in sets and $S^1$ at $t$ in
groupoids. The main family and the nonconstant nine-object example
preserve their left generation on old inputs; this example does not.
The two behaviors concern different hypotheses.

\begin{table}[htbp]
\centering
\begin{tabular}{@{}lcrrl@{}}
\toprule
Six-object family & $|D|$ & Sets & Groupoids & Generation on old inputs\\
\midrule
$C'_2$ of \cref{prop:compact} & 1 & 5 & 6 & Preserved\\
Cone family of \cref{prop:reversal} & 4 & 6 & 5 & Changed\\
\bottomrule
\end{tabular}
\caption{Sizes of least resolvers in two different six-object posets.
Preserving generation guarantees monotonicity under the stated comparison
hypotheses; changing generation can reverse the inclusion.}
\label{tab:directions}
\end{table}

\subsection{A resolution profile}

For a fixed ambient indexing category, starting level, and admissible family,
call the sequence $k\mapsto\Aup_k(D)$ its \emph{resolution profile},
where the least levels exist. The expression is a name for the comparison
being studied, not a replacement for established cofinality terminology.

The profile contains more information than a single result. For the growing
families it shows both the validity of a restricted resolution and the first
expressive level at which it ceases to suffice. The reversal profile instead
loses a formerly essential position when generation changes.
In \cref{thm:main}, the failure has a
specified location, namely $a$, a specified shape, a wedge of spheres,
and a specified detecting coefficient, an Eilenberg--Mac Lane space. What
would otherwise be an unspecified demand for ``more structure'' becomes a
calculable request.

This also clarifies a distinction between erasure and invisibility. The
nerve is a circle already in the first example. Set-valued coefficients do
not turn that nerve into a contractible space; they are unable to distinguish
it from a point by the maps under consideration. The distinction untested by the comparison is
therefore relative to that test. To say that a result preserves everything
relevant is to say something about what counts as a relevant test, not only
about the construction of the result.

\subsection{Why leastness is not enough}

A least object is uniquely determined up to the relevant equivalence once
the ordered collection of candidates and its predicate of adequacy are
fixed. Leastness can therefore express a strong form of internal necessity.
In our family no alternative full subposet is smaller and still meets the
specified test. The matching-support calculation excludes all such alternatives.

Consider an explicit inference, introduced here as a claim to be tested,
not attributed to any of the authors discussed: if an enrichment keeps the
generation of the diagrams and every earlier resolution comparison intact,
then it also keeps their least adequate resolving level intact. The premise
is stronger than agreement on selected examples. It preserves the entire
old domain of comparisons. Yet \cref{prop:persistence,thm:main} jointly
refute the conclusion, with arbitrarily long finite chains of different
least levels.

What survives is correctness on the earlier domain; what fails is its
sufficiency to determine the least adequate reconstruction on the enlarged
domain. This distinction is not a discrepancy between an informal reading
and a formal calculation. It occurs between formally stated properties,
with fixed indexing data and fixed generation. Leastness combines an order
of candidates with a predicate of adequacy. Persistence of old comparisons
preserves one part of the predicate's evidence without establishing the
completeness of its range.

For a reconstruction of dialectical necessity, the consequence is conditional
but exact. If its explanatory commitments give a reason to admit additional
tests, continued agreement on all earlier tests does not by itself secure
its earlier leastness claim. It must justify the original test domain,
prove a transport result, or accept a changed least resolution. If no such
reason is given, our theorem supplies a possible comparison rather than a
refutation. This is the distinction between a useful external analysis and
an internally demanded revision to which we now turn.

\section{Determinate preservation as an interpretive requirement}
\label{sec:preservation}

\subsection{A criterion with a limited demand}

We propose the following requirement for a mathematical reconstruction of
sublation:
\begin{quote}
State the relations claimed to represent the mediation and the distinctions
whose retention constitutes the proposed preservation; show how the
resolving construction retains them; and justify the treatment
of distinctions that become detectable under the reconstruction's own
licensed changes of language.
\end{quote}
This requirement of \emph{determinate preservation} is our methodological
proposal. It is not a quotation from Hegel or a theorem of category theory.
Its motivation is that preservation needs a determinate content if it is to
play an explanatory role. The theorem supplies an exact case in which an
unqualified preservation claim would conceal a choice.

The demand is intentionally weaker than invariance under every enrichment.
Not every new formal distinction deserves philosophical recognition. A set-based
model can be exactly appropriate for a question about ordinary
classification. A groupoid model may be appropriate when identifications
and automorphisms matter. The requirement is to defend this choice in
relation to the interpreted material, not to rank models by how much
machinery they contain.

A change of test becomes relevant in at least three distinguishable ways.
The reconstruction may itself use the additional structure to explain the
passage; its claimed applications may require that structure; or a close
reading may show that its original tests identify distinctions it needs to
retain. These are warrants that can be contested with textual and
mathematical evidence. Merely knowing that groupoids are available is not
such a warrant. Conversely, declaring that only sets are admissible does
not settle the issue if the explanatory argument already invokes coherent
identity. The present theorem diagnoses consequences once a warrant is
given; it does not manufacture that warrant.

Nor does the requirement insist that every revision add more objects.
In the generation-preserving families the least resolver grows;
\cref{prop:reversal} shows that a change of generation can instead
make it shrink. Another model might revise its equivalences, its morphisms,
or its admissible tests instead. The general lesson is about explicit
dependence on representational choices; the particular repairs in a
constructed poset do not constitute a universal recipe for dialectic.

\subsection{Retention without identity of status}

The philosophical attraction of sublation is that something can persist
without retaining its former independence. A categorical reconstruction can
give this thought more than one precise form. An object may occur as the
image of an embedding, as a retract, as a component of gluing data, or as
the value of a modality. These are different modes of dependence. Their
differences should be used rather than obscured by a common label.

Our resolution condition does one specific job: an earlier skeletal object
is recoverable through a later right extension. This captures a relation of
recoverability and changed presentation. It does not independently capture
every philosophical sense of negation, nor does it show that the earlier
concept has criticized itself. The model is informative because its
achievement can be stated exactly.

There is also a limit to the familiar picture of preservation as storage.
Keeping a copy of every earlier object need not preserve the relations
that made those objects what they were. In homotopy theory, retaining points
while suppressing paths can alter the character of the object. In an
inferential account, retaining sentences while changing which consequences
they license can alter their content. The analogy between these cases is
methodological, not an identification of paths with inferences. Both require
the interpretation to specify the relations its preservation concerns.

\subsection{What is established in the content, and what is available to us}

The distinction in Hegel's discussion of determinate being is especially
instructive. In online paragraph 193, a determination may be available to
our reflection before it is established as part of the concept's own
development. Such a comparison can elucidate or anticipate; that does not
give it the status of a ground of the transition \cite{HegelDasein}.

The author's running commentary, distinct from the earlier interactive
map, foregrounds this problem in its first two readings
\cite{AutioCommentary}. Its distinction between what is available
``for us'' and what has been posited within the development helps formulate
the interpretive demand. It is a statement of the problem, not independent
evidence that a particular reconstruction meets it.

Applied to this paper, the distinction asks why a coefficient enrichment
belongs to the represented content rather than only to our means of
investigating it. The answer cannot be that the richer formalism exists.
Nor can it be that the enrichment has produced an interesting theorem.
Those facts justify studying the comparison, but they do not determine
its interpretive status. A close reading must identify the relevant demand
of the content; a mathematical preservation theorem must then establish
what the chosen formal correlate achieves.

This yields an interpretation audit with a definite result. Measured
against the demands drawn from the becoming passage, our formal calculation
establishes recoverability and least repair only for its stipulated
mathematical data. It does not supply the moments'
self-transformation or explain their altered significance. For the status
of the coefficient category, it supplies a sensitivity result but no
immanent warrant for enrichment. The mathematical result is therefore
philosophically useful in an elucidatory role, while its promotion to a
derivation would require evidence it does not provide.

\subsection{Immanence and the status of the formal language}

Hegel's final discussion of method presents the relation between method and
content as internal rather than as the application of an independently
chosen technique \cite{HegelIdea}. A mathematical reconstruction necessarily
begins with some formal resources. It cannot escape that fact by describing
those resources as universal. It can, however, make their role an object
of the reconstruction rather than leaving it as an unexplained background.

Hentrup's account of rational proof supplies an independent philosophical
comparator. In section II, geometric demonstration can establish a
necessity belonging to its object while its progression remains the
prover's activity. In sections V--VI, rational proof requires mutual
mediation in which the terms' identities develop, rather than reciprocal
support between fixed terms \cite{Hentrup}. On this reading, the
difference between valid demonstration and adequate mediation is already
a substantive Hegelian problem. Our contribution is not to discover that
difference, but to make two dependencies of a formal reconstruction
calculable: the tests that determine its least resolver and the
relations whose interpretation remains to be justified.

The resolution profile offers one finite way to do so. The coefficient
category is first a parameter of the model; the profile then shows how
the model's verdict depends on it. This does not reproduce Hegel's claim
that the method determines itself. It provides a disciplined partial
analogue: the consequences of the choice of method are brought into the
argument and can be evaluated. A stronger claim of self-determination
would require additional work explaining why the sequence of admissible
tests is generated by the subject matter rather than chosen by us.

The gap should remain visible. Our theorem can force the next resolving
level \emph{if} the next coefficient category is admitted. It does not force
the admission of that category. Confusing these two steps would repeat,
at a higher level, the very substitution of external comparison for internal
development that the reconstruction is meant to illuminate.

\subsection{A reader audit: stable profiles and different relations}
\label{sec:readeraudit}

Coefficient variation is not the only choice a reconstruction makes.
The author's earlier interactive reader \cite{AutioReader} supplies a
concrete case in which to distinguish it from the choice of indexing
relations. The reader locates sections within larger sections and
opens their associated passages. Containment and succession are both
legitimate editorial relations. Neither, by itself, establishes the
dependence by which one determination generates another.

Consider two explicitly stipulated encodings of the same positions.
In one, an ancestor is below its descendants in a rooted tree. In the
other, the positions are totally ordered by a traversal compatible
with that ancestry. A synthetic book-header root is first in both.
It is an editorial device, not an identification with Pure Being;
likewise the second order records succession, not implication.

\begin{proposition}[A limit of the profile diagnostic]\label{prop:treechain}
Let $T$ be a finite rooted tree with root $b$, with at least two
children at every nonleaf vertex. Let $L$ be a total order on the same
vertex set extending the ancestor order, and write $t$ for its last
vertex. For $D=\{b\}$ and every $0\leq k\leq\infty$,
\begin{equation}\label{eq:treechain}
\Aup_k^T(D)=T,\qquad \Aup_k^L(D)=\{b,t\}.
\end{equation}
Thus both profiles are invariant under coefficient enrichment, although
their least subsets can differ by arbitrarily many objects.
\end{proposition}

\begin{proof}
At a leaf of $T$ the strict upper interval is empty. At a nonleaf it is
the disjoint union of the descendant subtrees rooted at its children.
Each component has an initial object, so its nerve is contractible;
there are at least two such components. Every vertex therefore belongs
to $E_k(T)$ at every $k$. In $L$, only $t$ has empty strict upper
interval; at every earlier vertex that interval has a least element.
Hence $E_k(L)=\{t\}$. Apply \cref{thm:intrinsic} and add the required
starting object $b$. Finite branching trees can have arbitrarily many
vertices.
\end{proof}

\begin{figure}[htbp]
\centering
\begin{tikzpicture}[x=1cm,y=.9cm,
 kept/.style={circle,draw=ink,fill=ink,inner sep=2.6pt},
 omitted/.style={circle,draw=ink,fill=white,inner sep=2.6pt},
 edge/.style={draw=ink!65,line width=.65pt},
 every node/.style={font=\small}]
\node[kept,label=left:{$b$}] (tb) at (0,0) {};
\node[kept,label=left:{$u$}] (tu) at (-.8,1) {};
\node[kept,label=right:{$v$}] (tv) at (.8,1) {};
\node[kept,label=above:{$w$}] (tw) at (-1.3,2) {};
\node[kept,label=above:{$z$}] (tz) at (-.3,2) {};
\foreach \s/\t in {tb/tu,tb/tv,tu/tw,tu/tz}
\draw[edge] (\s)--(\t);
\node[align=center] at (0,-.9) {Containment\\$\Aup_k^T(D)=T$};
\node[kept,label=right:{$b$}] (lb) at (4.3,0) {};
\node[omitted,label=right:{$u$}] (lu) at (4.3,.55) {};
\node[omitted,label=right:{$w$}] (lw) at (4.3,1.1) {};
\node[omitted,label=right:{$z$}] (lz) at (4.3,1.65) {};
\node[kept,label=right:{$v=t$}] (lv) at (4.3,2.2) {};
\foreach \s/\t in {lb/lu,lu/lw,lw/lz,lz/lv}
\draw[edge] (\s)--(\t);
\node[align=center] at (4.3,-.9) {Succession\\$\Aup_k^L(D)=\{b,t\}$};
\end{tikzpicture}
\caption{Two stipulated orders on the same positions. Filled vertices
are in the least resolver for every coefficient level. The invariance
of both profiles does not decide which relation, if either, models
the intended mediation. Arrows in these Hasse diagrams point upward.}
\label{fig:treechain}
\end{figure}

The active reader contains 120 positions in four displayed depths.
After excluding its one unnamed placeholder and adding the synthetic
root, there are 120 named-or-root positions. Every internal vertex
still has at least two children. The proposition therefore applies
to these proposed encodings. These counts describe the inspected
editorial data; they assert neither a canonical inventory of Hegelian
categories nor a mathematical structure derived from the text.

The contrast supplies a limit on our own diagnostic. Invariance of a
resolution profile is insufficient for interpretive adequacy. The
opposed leastness verdicts here are explained by branching and
terminality before any philosophical content is assigned to the
labels. This does not make either encoding defective: it specifies
what each achieves. In this case the unsettled choice concerns
relations; in the Boolean family the relations are fixed and the
varying choice concerns tests. The two choices require distinct
arguments.

A passage-grounded extension would therefore distinguish containment,
succession, and argued dependence, recording textual evidence for the
latter instead of silently replacing it with either of the former.
Composites and identifications between dependency paths would require
further justification before a category or higher category is claimed.
The present paper does not invent those relations in order to find
Hegelian loops. It identifies what such an extension would need to
establish before a homotopy calculation acquired that interpretation.

\subsection{A mathematical reading need not be a reduction}

The distinction between orientation and derivation has a direct textual
basis. In the final discussion of method, Hegel permits an advance survey
of the whole while denying that division into its already specified parts
establishes their transition (online paragraph 1813). The author's earlier
interactive reader \cite{AutioReader} makes this distinction concrete:
its hierarchy locates moments and opens the corresponding passages.
Such navigational success need not claim to establish their mediation.
Nor does a circular display by itself realize the return through mediation
described in paragraph 1814. A formal interpretation must say which of
these achievements it intends.

There is also a challenge to our own procedure. Paragraph 1812 rejects
the claim that criticism of cognition's instrument enjoys an automatic
priority over engagement with the subject matter: that criticism has
presuppositions of its own \cite{HegelIdea}. Determinate preservation
should therefore operate within the examination of a reconstruction,
rather than as a preliminary tribunal that settles its possibility in
advance. The test must answer to the content it claims to illuminate.
Our theorem identifies a dependence to be justified; it neither supplies
a presuppositionless standpoint nor decides that justification before
the interpretation has been attempted.

The value of a mathematical reading is not exhausted by translation into a
fixed formal calculus. It can reveal a question that was difficult to state,
separate claims previously run together, and supply counterexamples to an
overextended argument. Here the precise separation is between the
existence of a least resolver and the justification of its criterion of
resolution.

This supports a reciprocal rather than a reductive relation between
philosophy and mathematics. The philosophical concern with retention
motivates a comparison of coefficients; the mathematical comparison then
forces the philosophical concern to become more specific. Neither side
receives authority simply by borrowing the other's vocabulary. What matters
is whether the encounter changes the questions that can responsibly be
asked and the answers that can be defended.

\section{Contemporary alternatives and adjacent results}
\label{sec:literature}

\subsection{Progression, contradiction, and logical form}

Priest's model uses paraconsistent logic, but also specifies a concrete
retention condition: introducing the new predicate leaves the old
vocabulary's extensions and anti-extensions unchanged
\cite[section 3]{Priest}. In section 6 he expressly limits the aim to the
logical structure of progression, rather than an explanation of why a
progression occurs. These are precedents for specifying preservation and
scope, not errors our theorem corrects. His retention condition is
stipulated in the construction; ours is Kan-extension recoverability
with a proved leastness profile. Neither identification alone establishes
its adequacy to the text.

An adjunction does not assert $P\wedge\neg P$, and a homotopical
obstruction does not license or reject explosion. A useful comparison
must therefore identify shared explanatory tasks before treating
paraconsistency and modal adjunctions as competing answers to one
question.

Yang's published abstract distinguishes the accuracy of describing
dialectical progression from the ability to explain it \cite{Yang}.
This distinction is congenial to our separation of retention,
determination, and necessity. We do not rely on a detailed reconstruction
of Yang's argument, whose full text was not available for this study.
Our mathematical result provides an independent instance of the general
difference: a diagram can describe a proposed transition without
settling why its adequacy criterion is the appropriate one.

Kolman's account treats Hegel's relation to logical form through the
reconsideration of fixed determinations rather than the mere addition of
one more logical form \cite{Kolman}. Our proposal shares that concern
with the status of a formal standpoint, while giving it a different
technical expression. Changing coefficients does not model a history of
logic. It gives a controlled comparison in which what a formalism can
detect becomes part of the account of its verdict.

Brown offers a positive example of making preservation explicit.
Definition 2.6 calls a morphism $f:X\to Y$ of incompatibility frames
conservative when every finite multiset $\Gamma$ is incompatible in $X$
exactly when $f(\Gamma)$ is incompatible in $Y$. Proposition 2.1 proves
this property for the adjunction unit; the signed construction inherits
it in section 3.2 \cite{Brown}. The universal construction thus has a
separately stated preservation achievement. It differs from our
Kan-extension recoverability. The methodological point is that specifying
an adjunction and specifying what it retains are distinguishable tasks.

\subsection{What recent Aufhebung calculations already achieve}

The work of Kennett, Riehl, Roy, and Zaks computes actual skeletal and
coskeletal bounds in ordinary presheaf topoi \cite{KennettEtAl}. Their
results show why an unspecified ``next dimension'' is not an adequate
description of Aufhebung. The precise bound depends on the indexing
category, and low-dimensional cases require attention.

The recent solution by Hora, Kamio, and Maehara to Lawvere's fourth
problem computes the relation for symmetric simplicial sets
\cite{HoraEtAl}. This removes that calculation from the list of open
targets. Our family changes a different parameter: its indexing category
is fixed while its coefficients vary. The elementary topology involved
does not make it a replacement for the substantially deeper combinatorics
of their theorem.

Menni's work on the least subtopos containing the discrete skeleton of
$\Omega$ shows the power of replacing a large containment condition by
a distinguished generator \cite{Menni}. It also makes essentiality a
substantive issue. One cannot freely identify a least subtopos with a
least essential resolving level. This is why our higher theorem names
its admissible full-subcategory levels explicitly instead of relying on
an unstated lattice classification.

Dey--Lesnick and Br\"ustle--Desrochers--Leblanc address the least
full subposets needed for ordinary limit or colimit computation
\cite{DeyLesnick,BrustleEtAl}. They are direct antecedents of the
$k=0$ upper-interval mechanism, while their treatment of nonfull
relation-sparse subposets optimizes a different admissible class. Our
calculation combines a truncated version of this mechanism with explicit
sphere obstructions to compare coefficients. It does not reopen the
ordinary least-full-subposet problem as if that literature did not exist.

\subsection{What a checked construction preserves}

DeBruce's Cubical Agda tutorial offers a useful contemporary example
because it places explicit limitations beside its formal declarations.
Chapter 4 introduces a higher inductive type \texttt{Becoming} with two
points, \texttt{being} and \texttt{nothing}, and a path between them
\cite{DeBruce}. These points are new inhabitants of a common carrier.
They are not the unit and empty types used for Being and Nothing in
Chapter 2. The distinction matters: a path between those earlier types
would transport an inhabitant of the unit type into the empty type.
The consistent construction therefore depends on changing what is
related, not simply on replacing judgmental equality with path equality.
An interpretation must explain which roles of the earlier assignments
survive this change.

The same chapter declares an \texttt{Aufhebung} record containing a
higher unity and a function from the lower moment to the higher. It
explicitly leaves an instance unconstructed and acknowledges missing
compatibility requirements; it supplies no leastness theorem. These are
useful limits on the achievement, not reasons to dismiss it. The tutorial
makes an interpretive audit concrete: naming, providing a structure,
retaining a specified property, and proving minimality are distinguishable
accomplishments. We assess the inspected declarations and their
interpretation; we have not independently rerun that project's Agda build.

\subsection{Higher modal semantics is infrastructure, not a transfer theorem}

Uemura's internal languages for diagrams of infinity-logoses connect
modalities, gluing, and higher categorical semantics \cite{Uemura}.
Shulman's multimodal adjoint type theory provides a systematic treatment
of modal adjunctions and their semantics \cite{Shulman}. Myers and Riley
study multiple commuting cohesions \cite{MyersRiley}. These works make
the higher setting substantially more precise than an informal instruction
to replace sets by types.

The computational side is also important. Gratzer's normalization result
relates type checking to equality in the chosen mode theory
\cite{Gratzer}; the combination of cubical and multimodal type theory
provides another route toward expressing homotopical modal arguments
\cite{AagaardEtAl}. None of these metatheorems selects the mode theory
appropriate to Hegel, and none by itself says that a given set-valued
resolution has the same least solution in spaces. They supply languages
in which such questions can be asked rigorously.

Protin's recent discussion of Hegel and topology offers mathematical
illustrations and hermeneutic proposals \cite{Protin}. The contrast with
our purpose is one of task: we isolate a preservation claim and calculate
its failure of transport. An illustration can illuminate a passage without
claiming to establish a universal property; a universal property can be
proved without establishing that the passage has been interpreted correctly.
The distinction protects both kinds of work from exaggerated expectations.

\section{Conclusion: the conditions of a necessary transition}
\label{sec:conclusion}

The finite families give definite mathematical answers. A fixed initial
level can have successively different least resolutions when its admissible
coefficients range from sets through higher truncated spaces. The entire
sequence is calculable; its obstructions are explicit; its leastness holds
against every full-subcategory competitor. The earlier answer remains
true within its domain, while a richer domain detects what it did not test.
When generation changes as well, the reversal family shows that the least
resolver can instead shrink. Preserving the earlier generation is therefore
a substantive condition on a comparison, not a harmless convention.

The reader audit supplies the converse restraint. Containment and
succession, encoded on the same positions, can yield sharply different
least resolvers while both profiles remain coefficient-invariant.
Robustness under a change of tests is therefore no certificate that the
relations being tested express the intended mediation. Together the
examples separate two tasks: accounting for the admitted relations,
and accounting for the domain over which their adequacy is tested.

For a mathematical reading of Hegel, the result recommends neither a
retreat from formalization nor an unqualified preference for higher
structures. It recommends a more demanding account of preservation.
The interpretation must identify the content of what survives, the
equivalences under which it survives, and the tests by which its survival
is judged. As the interpretation proceeds, those choices must answer
to the text and explanatory task. A universal property can then help
explain the conditional necessity of its result without being asked
to justify the representation by itself.

Lawvere's approach is especially valuable because it permits this demand
to become mathematics. A proposed resolution can fail, and the failure
can have a shape. In our example that shape is a sphere invisible to one
coefficient category and visible to the next. The philosophical lesson
is not that every hidden sphere is a dialectical moment. It is that a
formal account of a necessary transition owes an account of what its
notion of necessity can see.
